\chapter{引言}
\section{选题背景及意义}
弹性波检测技术[]指的是在不损坏被检测对象的情况下，
对检测对象输入弹性冲击波，通过分析接收到的波速分布图像来检测结构内部缺陷和结构完整性。
该方法在土木工程领域结构损伤检测中得到了广泛的应用，
如桥梁、隧道等混凝土结构损伤监测[]，地震、石油勘探[]，建筑结构评估[]等问题中。
弹性波检测技术中弹性波在结构中的传播可采用弹性波偏微分方程进行描述和分析，
该问题也称为弹性波问题[]。通过弹性波问题分析波在介质中的传播过程可确定最优的弹性波类型、
优化检测位置等，从而提高检测精度。然而，对于复杂求解区域，
弹性波问题难以进行解析求解，通常需要采用有限元[]等数值近似方法模拟波在介质中的传播过程。

在传统有限元弹性波分析[]中，空间域和时间域通常分开进行近似离散，
从空间域进行离散的方法属于边值问题，即所有边界条件均是已知的，包括位移边界条件、
应力边界条件等。该问题通常采用有限拓扑单元进行离散，适用于复杂不规则求解区域，
如含裂纹、空洞的不规则域。求解过程可采用伽辽金法进行整体求解，
整体求解过程计算精度高，可高度并行化，计算效率高。
以二维问题为例，三角形单元与四边形单元的混合使用可有效处理曲率边界，
而三维问题中四面体、六面体及多面体单元的灵活组合，
能够精准拟合地质界面或工程结构的复杂轮廓。值得注意的是，
空间域离散本质上是将连续的波动方程转化为分片多项式插值的变分问题，
其数学基础源于加权余量法中的伽辽金原理 —— 通过选择与形函数同空间的权函数，
构建单元刚度矩阵与荷载向量，进而通过组装形成全局离散方程组。

时间区域属于初值问题，即只已知初始时刻位移与速度，时间末端边界条件未知。
该问题通常采用迭代方法对时域进行积分，迭代过程中误差会随着迭代步数的增加而累计，
时间末端的计算精度低。并且，迭代步计算过程需要采用上一迭代步的计算结果，
无法进行高度并行化。与此同时，时间迭代步的步长也将影响计算的稳定性[]，
过大的步长将导致计算结果不符合精度要求。
通过对迭代步计算结果的后误差估计[]可确定下一步的时间步长，然而，
空间区域的网格畸变问题[]会导致精度下降，使得保证精度的时间步过小，计算终止。
另一类迭代算法则是通过哈密顿正则方程生成的保辛算法(Symplectic integrators)[]，
该类方法由冯康系统地总结了先前算法提出[]。在拉格朗日力学中，
动力问题需满足变分原理或称哈密顿原理(Hamilton’s Principle)，
通过对变分原理得到的欧拉-拉格朗日方程(Euler-Lagrangian equation)进行正则变换，
可得等价得哈密顿正则方程。采用传统的迭代方法求解哈密顿方程时可确保能量、
角动量守恒，即满足保辛性(symplecticity)[]。目前，大多数保辛算法用于分子动力学、
天体物理、量子力学等领域[]，如蛙跳法[]、Verlet法[]，隐式中点法[]、
半隐式欧拉法[24]、分区Runger-Kutta法[]等。在结构动力学领域，Simo等人[]、
Kane[]分别验证了Crank-Nicholson法、Newmark- 法可满足保辛性。

时空混合离散方法[]是求解波动问题另一种思路，
该方法通过哈密顿原理将时间区域和空间区域同时采用有限元形函数进行离散，
时间区域和空间区域均采用变分原理进行求解。相较于主流迭代算法，
时空混合离散无需在时间步进行迭代，可同时求解所有时刻、所有空间区域节点系数。
求解过程可高度并行化，计算效率高。同时，时空混合离散网格局部细化过程易于实现，
空间区域的局部细化将不会对后续时间区域构成影响。然而，
哈密顿原理要求时间末端虚位移必须为零，但往往该问题时间末端位移并不知道，
直接施加虚位移节点系数为零将导致刚度矩阵不为方阵[]。
解决该问题最直接方式是将将初始时间的强制边界条件施加到时间末端，
为此保证整个系数矩阵是方阵，从而可以进行求解[]。这种做法存在明显的局限性，
该方法无法对网格进行任意离散，必须保证初始时刻的节点和最后时刻的节点保持一致。
此时，时空离散网格难以进行局部加密，限制了其在复杂问题中的应用。
为此，Hughes等人提出了一种间断伽辽金法[]，
该方法在泛函变分中添加非连续的边界条件使弱形式满足变分一致性，
此时弱形式无需要求虚位移在时间末端为零。
但为添加的非连续边界条件将导致弱形式的正定性无法得到满足，
需要添加额外的稳定项来保证计算精度。同时，间断伽辽金方法仍然无法实现绝对非均布的离散，
只能在均布的时域块内做非均布计算，对于具有高度局部特性的问题时仍不够理想。
Wriggers和Junker采用虚单元法构造非协调单元进行时间区域与空间区域的整体离散[]，
也是属于间断伽辽金范畴，该方法可以使用高阶单元来提高计算精度，但同样需要添加稳定项。
采用配点型方法也可以规避处理时间末端边界条件的问题[]，
并引入正交基函数离散哈密顿正则方程进行求解。该方法求解精度与配点位置相关，
为保证计算精度也需要采用均布网格进行离散。目前，
时空混合离散方案均无法做到绝对的时空混合离散，限制了该类方法在实际工程中的应用和推广。

综上所述，有限元模拟弹性波检测的分析方法经历了从传统迭代方法到基于哈密顿原理的时空混合离散方法的演变。
这些方法的发展不仅提高了分析的效率和精度，还增强了模拟的稳定性。
但哈密顿原理在完全时空离散有限元法处理弹性波问题仍然存在这时间末端边界和无法任意使用局部加密等问题，
这些问题成为发展哈密顿原理计算动力系统数值分析方法的瓶颈。
本论文将系统研究完全时空离散有限法元计算动力系统问题，
探究哈密顿变分原理作有限元分析时的边界问题，从而确定处理时间末端边界问题的必要条件。
同时，建立一个在有限元伽辽金数值计算中能完全时空混合离散的方案。
从而得到更高效、精确的数值分析方法，从而能稳定应用到解决弹性波问题当中，
为有限元模拟弹性波检测技术应用提供可靠精确的数值分析工具。
\section{国内外研究历史及现状}

\section{本文主要内容}